\section{Microservice implementation: Search and deployment}
\label{sec:search}

\subsection{Search Service}

Search is written in Java 17 with Javalin and Gson, and implements \code{api/openapi.yaml}. It never writes and
mounts no volume: everything it needs is in MongoDB, which makes it easy to move to another machine in Stage 3.
Its main classes are:

\begin{table}[H]
\centering\small
\begin{tabularx}{\linewidth}{@{}L{4.3cm}X@{}}
\toprule
\textbf{Class} & \textbf{Responsibility} \\
\midrule
\code{SearchApi} & Validates the parameters and writes the HTTP answers \\
\code{SearchEngine} & Runs one query on one snapshot \\
\code{IndexSnapshotLoader} & Builds a snapshot of the published version from MongoDB \\
\code{IndexSnapshot} & Immutable: version, published ids and their metadata \\
\code{SnapshotHolder} & Keeps the current snapshot in an \code{AtomicReference} \\
\code{SnapshotRefresher} & Checks \code{CURRENT} periodically and on demand, on one thread \\
\code{PublicationEvents} & Listens to \code{IndexPublished} in the background \\
\code{MongoPostingSource} & Reads the posting lists of the terms with one query \\
\bottomrule
\end{tabularx}
\caption{Structure of the Search Service.}
\end{table}

\textbf{One query.} \code{SearchEngine} reads the posting lists of the terms. It walks the shortest list and
keeps an id only if every other list has it, using binary search, so the cost depends on the rarest term. Then
it drops the ids that the snapshot does not publish, applies the filters on the metadata in memory, orders all
the matches and cuts the requested page.

\textbf{Consistency.} Indexing may already have written books that are not published, so Search must not show
everything it finds. Each request takes the snapshot once when it starts and works only with that one, so it
cannot mix versions, and it reports that version as \code{index\_version}. The loader reads \code{CURRENT} first
and then the books up to that version. A book keeps the version in which it first became visible, so whatever
moment the loader runs, it sees all of a version and never part of the next one. Replacing the snapshot is one
atomic write: no request waits for it, and no thread can see one half built.

\textbf{Reloading.} A new version reaches Search in two ways. The \code{IndexPublished} event makes it react at
once, but it is only a wake-up: its content is not used, and whoever receives it reads \code{CURRENT} from
MongoDB. The check of \code{CURRENT} every \code{SEARCH\_POLL\_SECONDS} (5 s) is the backup for lost events.
Both run the reload on the same thread, so a slow reload never blocks the Hazelcast thread and two
reloads never overlap. If a reload fails, the replica keeps serving the snapshot it has.

A failure found in testing shows why both are needed. The ringbuffer of a new Hazelcast cluster starts again at
sequence 0, while the subscription keeps waiting at the sequence it had reached. After restarting both members,
Search received no event until as many versions had been published, and only the periodic check kept it up to
date. Each reconnection now subscribes again and checks \code{CURRENT} once; a test restarts the cluster
repeatedly and publishes one event at a time, as Indexing does.

\textbf{Availability.} Search starts even if MongoDB or Hazelcast are down, or nothing has been published.
\code{/status} answers from the first moment, because connecting to Hazelcast happens in the background after the
HTTP server starts. Searches answer 503 until a version is loaded. If Hazelcast is down, Search keeps working with
the periodic check; if MongoDB is down, searches answer 503 \code{STORAGE\_UNAVAILABLE} and recover by themselves.

\textbf{Concurrency.} A request runs on a thread of the Jetty pool, whose maximum is explicit in
\code{SEARCH\_HTTP\_THREADS}: 250, the default value of Javalin 6.1.3, with a minimum of 8 and a 60~s idle
timeout. These threads share only immutable data, so they need no lock. One \code{MongoClient} per process has a pool
of \code{SEARCH\_MONGO\_POOL\_SIZE} connections (50), and a request that waits more than 5~s for one answers
503 instead of keeping its thread busy for the 2 minutes of the driver's default.

\textbf{Browsers.} \code{CORS\_ORIGINS} lists the origins allowed to call the API from a web page; an invalid value
stops the start.

\subsection{Deployment}
\label{sec:deployment}

\code{docker compose up} starts everything on one host from \code{docker-compose.yml} and \code{.env}, which
holds every default. All Java services use one multi-stage template, \code{docker/service.Dockerfile}, with the
Maven module as an argument.

\begin{table}[H]
\centering\small
\begin{tabularx}{\linewidth}{@{}L{2.5cm}L{1.3cm}L{3.1cm}L{1.9cm}X@{}}
\toprule
\textbf{Service} & \textbf{Replicas} & \textbf{Published on the host} & \textbf{Datalake} & \textbf{Limits} \\
\midrule
\code{nginx} & 1 & 8080 & none & 0.5 CPU, 128 MiB \\
\code{search} & 1--4 & none & none & 2 CPU, 1 GiB each \\
\code{indexing} & 1 & none & read-only & 2 CPU, 1 GiB \\
\code{ingestion} & 1--4 & none & read-write & 1 CPU, 512 MiB each \\
\code{control} & 1 & 7000, local only & read-only & 0.5 CPU, 384 MiB \\
\code{hz-1}, \code{hz-2} & 2 & 5701, 5702, local only & none & 1 CPU, 768 MiB each \\
\code{mongo} & 1 & 27017, local only & none & 2 CPU, 1 GiB \\
\bottomrule
\end{tabularx}
\caption{Containers of the deployment. One replica of each adds up to about 5.7~GB of memory limits.}
\end{table}

\begin{itemize}
  \item \textbf{Replicas.} \code{docker compose up -d --scale search=3} is enough: nginx
        (\code{scripts/nginx.conf}) does round-robin over the replicas, which it finds with the DNS of Docker and
        looks up again while running.
  \item \textbf{Failures.} The public API is read-only, so nginx can safely repeat a request on another replica if
        one dies. It does not do so on a 503: a replica that answers is alive, and the others would answer the same.
        Killing one of 3 replicas during a test lost none of 2,000 requests.
  \item \textbf{Permissions.} Only Ingestion can write in the datalake, and Docker enforces it. Hazelcast has no
        authentication, so only nginx listens on all the interfaces.
  \item \textbf{Fixed resources.} Every container has limits, and the heap is a fixed share of its memory, so the
        benchmarks can be repeated and report what they used.
  \item \textbf{Monitoring.} \code{scripts/monitor.sh} records the CPU and memory of each container and the load of
        the host as CSV. A sample of the containers takes about 2~s, because of \code{docker stats}.
\end{itemize}

\textbf{Replicas and threads are different things.} A thread shares its process with the others, including the
snapshot. A replica is another process with its own memory, its own snapshot and its own connection pool, reached
through nginx. Both share the CPU, memory and disk of the host, and all the replicas share one MongoDB.
